\section{Architecture Overview}
AeroCycle is organized as two cooperating layers on each node, as shown in
Figure~\ref{fig:arch}: an application-layer \emph{Volatility Estimator} that consumes raw
sensor readings and produces a target sleep interval, and a MAC-layer \emph{Negotiation
Module} that reconciles this target against link-quality feedback before committing to a
sleep decision.

\begin{figure}[H]
\centering
\begin{tikzpicture}[
  block/.style={draw=nustgreen,thick,rounded corners=2pt,fill=nustmist,minimum width=3.4cm,minimum height=1.1cm,align=center,font=\small},
  arr/.style={-latex,thick,nustslate}
]
\node[block] (sensor) {Sensor\\Sampling};
\node[block,right=1.4cm of sensor] (vol) {Volatility\\Estimator};
\node[block,right=1.4cm of vol,fill=nustgold!25,draw=nustgold!70!black] (neg) {Negotiation\\Module};
\node[block,below=1.1cm of neg] (radio) {LoRa\\Radio};

\draw[arr] (sensor) -- (vol);
\draw[arr] (vol) -- node[above,font=\scriptsize]{target interval} (neg);
\draw[arr] (neg) -- (radio);
\draw[arr] (radio.west) -- ++(-3.9,0) |- node[left,pos=0.75,font=\scriptsize]{link SNR} (neg.south);
\end{tikzpicture}
\caption{AeroCycle two-layer architecture: application-layer volatility estimation feeds a MAC-layer negotiation module.}
\label{fig:arch}
\end{figure}
